(:term:base.terms.long_call:) is a (:term:base.terms.segment_changing_control_transfer:). It first constructs the (:term:base.terms.long_return_frame:) on SS:SP, then commits the new CS and PC together. The two-operand encoding takes the new CS image from an (:term:base.terms.rn_register:). Its target EA has Q interpretation width: a memory operand reads a 64-bit offset and a register or immediate operand supplies its value directly.

\begin{enumerate}
\item Evaluate the target operand completely in the old architectural context.
\item Validate the complete new CS image.
\item Validate the target offset against the new CS and probe executable translation at the resulting (:term:base.terms.linear_address:).
\item Probe both 8-byte long-return-frame stores through the old SS context.
\item Commit the two frame stores, decremented SP, new CS, and new PC as one successful control-transfer operation.
\end{enumerate}

\textbf{Fault atomicity.} No frame word, SP, CS, or PC update is visible when any prior step faults.

The (:term:base.terms.long_return_frame:) is 16 bytes in the old SS context: the (:term:base.terms.continuation:) PC occupies offset 0 and the return CS image
occupies offset 8. An invalid new CS image raises (:ref:base.events.EXCEPTION.INVALID_CONTROL_IMAGE:); target translation or either
frame-store translation failure raises the applicable \texttt{ADDRESS\_TRANSLATION} family event.

\begin{enumerate}
\item Evaluate the target operand completely in the old architectural context.
\item Validate the complete new CS image.
\item Validate the target offset against the new CS and probe executable translation at the resulting (:term:base.terms.linear_address:).
\item Probe both 8-byte long-return-frame stores through the old SS context.
\item Commit the two frame stores, decremented SP, new CS, and new PC as one successful control-transfer operation.
\end{enumerate}

\textbf{Fault atomicity.} No frame word, SP, CS, or PC update is visible when any prior step faults.

The (:term:base.terms.long_return_frame:) is 16 bytes in the old SS context: the (:term:base.terms.continuation:) PC occupies offset 0 and the return CS image
occupies offset 8. An invalid new CS image raises (:ref:base.events.EXCEPTION.INVALID_CONTROL_IMAGE:); target translation or either
frame-store translation failure raises the applicable \texttt{ADDRESS\_TRANSLATION} family event.
